\documentclass{beamer}

% NPBT ("Norman's Pandoc Beamer Themes") ships several colour variants;
% NPBT_FOM is used here rather than the bundled NPBT_SC demo variant, whose
% "line" graphic is shipped under upstream's own tree as
% lNPBT_SC_linie.png (stray leading "l") while the theme code composes the
% path as NPBT_SC_linie -- a filename typo in the upstream package that
% throws "File not found" under any colour scheme selecting NPBT_SC. The
% FOM colour scheme's assets are named consistently and is used instead.
\usetheme{NPBT}
\usecolortheme{NPBT_FOM}

\usepackage[utf8]{inputenc}
\usepackage[T1]{fontenc}
\usepackage{graphicx}
\usepackage{booktabs}
\usepackage{tikz}

\title{Field Notes from the Boreal Survey}
\subtitle{A Beamer starter built on the NPBT theme}
\author{Alex Rivera}
\institute{Department of Environmental Science}
\date{\today}

\begin{document}

\begin{frame}
  \titlepage
\end{frame}

\begin{frame}{Outline}
  \tableofcontents
\end{frame}

\section{Introduction}

\begin{frame}{Why This Survey}
  \begin{itemize}
    \item Boreal forest cover has changed measurably over the last decade
    \item Existing satellite estimates disagree by up to 12\% in mixed stands
    \item A ground-truthed transect survey closes that gap
    \item This deck walks through method, sites, and early results
  \end{itemize}
\end{frame}

\section{Method}

\begin{frame}{Survey Design}
  \begin{columns}[T]
    \begin{column}{0.5\textwidth}
      \textbf{Sampling}
      \begin{itemize}
        \item 48 transects, 500 m each
        \item Stratified by canopy density
        \item Repeated across two seasons
      \end{itemize}
    \end{column}
    \begin{column}{0.5\textwidth}
      \textbf{Instruments}
      \begin{itemize}
        \item Handheld LiDAR
        \item Multispectral drone passes
        \item Soil moisture probes
      \end{itemize}
    \end{column}
  \end{columns}
\end{frame}

\begin{frame}{Site Coverage}
  \begin{table}
    \centering
    \begin{tabular}{lrr}
      \toprule
      Region & Transects & Area (km$^2$) \\
      \midrule
      Northern block & 19 & 41.2 \\
      Central block & 17 & 33.6 \\
      Southern block & 12 & 22.9 \\
      \bottomrule
    \end{tabular}
    \caption{Transect allocation by region.}
  \end{table}
\end{frame}

\section{Results}

\begin{frame}{Canopy Density Estimates}
  \begin{itemize}
    \item Ground survey and satellite estimates converge in dense stands
    \item Divergence concentrated in mixed and edge stands
    \item Drone multispectral passes explain most of the gap
  \end{itemize}
  \vfill
  \begin{block}{Headline result}
    Ground-truthed density is within 3\% of satellite estimates once edge
    stands are excluded, down from a 12\% gap in the raw comparison.
  \end{block}
\end{frame}

\begin{frame}{Seasonal Comparison}
  \begin{figure}
    \centering
    \begin{tikzpicture}
      \draw[->] (0,0) -- (6,0) node[right] {\small season};
      \draw[->] (0,0) -- (0,3.2) node[above] {\small density};
      \draw[thick] (0.3,0.6) -- (1.6,1.4) -- (2.9,1.1) -- (4.2,2.3) -- (5.5,2.6);
      \foreach \x/\y in {0.3/0.6, 1.6/1.4, 2.9/1.1, 4.2/2.3, 5.5/2.6}
        \fill (\x,\y) circle (2pt);
    \end{tikzpicture}
    \caption{Measured canopy density across five survey passes.}
  \end{figure}
\end{frame}

\section{Conclusion}

\begin{frame}{Summary}
  \begin{itemize}
    \item Ground-truthed transects reduce estimate error to under 3\%
    \item Edge stands remain the dominant source of disagreement
    \item Next season adds 15 transects targeting edge zones specifically
  \end{itemize}
\end{frame}

\begin{frame}{Thank You}
  \centering
  \Large Questions?\\[1em]
  \normalsize alex.rivera@example.edu
\end{frame}

\end{document}
